\documentclass[aspectratio=169]{beamer}
\usepackage[english]{babel}
\usepackage[utf8]{inputenc}
\usepackage[T1]{fontenc}
\usepackage{graphicx}
\usepackage{booktabs}
\usepackage{xcolor}

\usetheme{metropolis}
\setbeamercolor{background canvas}{bg=white}
\usecolortheme{default}
\setbeamertemplate{navigation symbols}{}

\setbeamertemplate{footline}
{
  \leavevmode%
  \hbox to \paperwidth{%
    \hfill
    \begin{beamercolorbox}[ht=3ex,dp=2ex,right]{}
       \color{black}
       \insertframenumber{} / \inserttotalframenumberhspace*{3ex}
    \end{beamercolorbox}%
  }%
  \vskip0pt%
}

\title{Modernizing JADE: Autonomous Agentic Pipeline}
\subtitle{M1 Checkpoint: Self-Improving Skills for Codebase Evolution}
\author{Mateusz Jarosz, Sebastian Rydz}
\institute{Warsaw University of Technology}
\date{\today}

\begin{document}

\begin{frame}
  \titlepage
\end{frame}

\begin{frame}{Project Goals}
\begin{itemize}
  \item \textbf{Goal:} Migrate the legacy JADE framework (Java 1.5) to modern LTS versions (up to 21) without human-in-the-loop during execution.
  \item \textbf{Innovation:} Move beyond standard LLM chat by building a system that \textbf{generates its own migration tooling}.
  \item It identifies failure patterns, creates reusable skills, and compounds knowledge across version jumps.
\end{itemize}
\end{frame}

\begin{frame}{System Structure}
\begin{itemize}
  \item \textbf{The Orchestrator:} Claude Code acts as the central router, executing specific tasks via modular files.
  \item \textbf{The 6-Step Loop:}
  \begin{itemize}
    \item \textbf{Scan} (detect idioms), \textbf{Plan} (order execution), \textbf{Copy} (freeze baseline)
    \item \textbf{Apply} (run tools), \textbf{Verify} (benchmark), \textbf{Capture} (feed failures to Skill-Creator)
  \end{itemize}
  \item \textbf{Skill Architecture:} Differentiates between Hand-authored skills (e.g., \texttt{jade-phase0-scanner}) and Auto-generated skills stored in a versioned registry.
\end{itemize}
\end{frame}

\begin{frame}{Conceptual Changes \& Industry Validation}
\begin{itemize}
  \item \textbf{The Pivot:} Shifted from ``teaching an agent to code'' to ``equipping an agent with deterministic enterprise skills''.
  \item \textbf{AWS Summit Inspiration:} Attended the ``Updating Legacy Applications'' panel and consulted with AWS engineers.
  \item \textbf{Validation:} Discovered the \textbf{AWS Transform CLI} uses a highly similar approach.
  \item Analyzed their open-source repository (\texttt{aws-transform-custom-samples}) and adapted their \texttt{codebase-analysis} and \texttt{java-modernization} blueprints into our custom Claude skills.
\end{itemize}
\end{frame}

\begin{frame}{Working Components \& Progress}
\begin{itemize}
  \item \textbf{Component 1: The Scanner (\texttt{jade-phase0-scanner}).} A grep-based idiom detector that outputs structured severity flags (e.g., \texttt{RAW\_INST\_FILES}, \texttt{BLOCKER}).
  \item \textbf{Component 2: Migration Skills.} Built \texttt{jade-1.5-to-1.6-raw-types} and \texttt{jade-1.5-to-1.6-enhanced-for} to handle specific, isolated refactoring tasks safely.
  \item \textbf{Component 3: Verification Sandbox.} Implemented \texttt{run-benchmark.sh} to measure unchecked warning deltas and enforce compile-success gates.
\end{itemize}
\end{frame}

\begin{frame}{Pathway Going Forward}
\begin{itemize}
  \item \textbf{Late-May:} Fully automate the \textbf{Skill-Creator} (transitioning from manually capturing failure patterns to autonomous generation).
  \item \textbf{Early-June:} Tackle the complex 1.7 $\rightarrow$ 1.8 LTS paradigm shift (Streams/Lambdas).
  \item \textbf{Mid-June:} Execute the 1.8 $\rightarrow$ 11 LTS jump (handling the removal of CORBA/IIOP dependencies).
\end{itemize}
\end{frame}

\begin{frame}[plain]
  \centering
  \Huge \textbf{Thank You!}

  \vspace{1.2em}
  \Large \textbf{Questions?}
\end{frame}

\end{document}